% ICIP 2026 submission - uses spconf.sty + IEEEbib.bst
\documentclass{article}
\usepackage{spconf,amsmath,amssymb,graphicx,booktabs}
\usepackage[hidelinks]{hyperref}
\usepackage{enumitem}
\usepackage{microtype}

\title{RFI DETECTION IN SENTINEL-1 SAR IMAGERY VIA PSEUDO-LABEL
DISTILLATION AND MULTI-RESOLUTION ENSEMBLE}

\name{Anas Elghoudane}
\address{Independent Researcher\\
\texttt{anaselghoudane@gmail.com}}

\begin{document}
\ninept
\maketitle

\begin{abstract}
We present a detection pipeline for Radio Frequency Interference (RFI) in
Sentinel-1 Synthetic Aperture Radar (SAR) quicklook imagery, achieving
0.4776 mAP$_{50\text{-}95}$ on the held-out test set of the ESA ClearSAR
Track 1 challenge. RFI manifests as thin horizontal stripes (median
height $\approx$10\,px, median width $\approx$140\,px) in variable-size
Sentinel-1 RGB quicklooks (median $\approx$515$\times$342\,px), a
small-object geometry that differs markedly from standard COCO
categories. The pipeline combines (1) a YOLO26x base detector with
SAR-specific augmentation and pixel-space Normalised Wasserstein
Distance (NWD) regression loss; (2) pseudo-label distillation from a
multi-fold ensemble teacher onto unlabelled test imagery; (3)
multi-resolution test-time augmentation (TTA) fused via Weighted Box
Fusion (WBF); and (4) selective cross-fold ensembling that discards
the domain-overlapping fold before final fusion. All ablation numbers
are computed on a 401-image leak-isolated meta-validation set using
pycocotools. Pseudo-distillation contributes $+0.041$
mAP$_{50\text{-}95}$ over the per-fold baseline; multi-resolution TTA
adds $+0.034$; selective fold dropping adds $+0.002$, for a total
$+0.077$ over the per-fold baseline.
\end{abstract}

\begin{keywords}
Ensemble, pseudo-label distillation, radio-frequency
interference, SAR imagery, small object detection.
\end{keywords}

\section{INTRODUCTION}
\label{sec:intro}

Synthetic Aperture Radar (SAR) sensors aboard Sentinel-1 provide
continuous global monitoring at C-band (5.4\,GHz). As terrestrial wireless
infrastructure has expanded, Radio Frequency Interference from
communication networks and other emitters has become an increasingly
common data quality problem \cite{njoku2005amsre}. RFI
corrupts focused SAR imagery as bright horizontal azimuth lines and, in
quicklook products, as rectangular stripes spanning the full image width
at affected azimuth positions.

The ESA ClearSAR Track 1 challenge \cite{clearsar2026} frames RFI
detection as a standard object detection task: given a variable-size
(median $\approx$515$\times$342\,px) 8-bit RGB quicklook, localise all
RFI-contaminated azimuth lines
with axis-aligned bounding boxes and report results in COCO format
\cite{lin2014coco}. The metric is COCO mAP averaged over IoU thresholds
0.50--0.95 (mAP$_{50\text{-}95}$).

RFI stripes are geometrically unlike any COCO category: a median
per-box aspect ratio of $\approx$8:1 (median height $\approx$10\,px,
median width $\approx$140\,px), with box heights spanning roughly
4--120\,px (5th--95th percentile) and widths up to $\approx$500\,px.
This geometry creates several detection challenges:
\begin{itemize}[noitemsep,leftmargin=*]
  \item IoU-based regression loss provides no gradient when predicted
  and target boxes do not overlap --- common when height predictions are
  off by only a few pixels.
  \item SAR side-looking geometry breaks horizontal symmetry: left--right
  reflections change the physically meaningful look direction.
  \item Multi-scale test-time augmentation above 640\,px amplifies
  speckle noise relative to the signal, hurting recall at larger
  resolutions.
\end{itemize}

We address these challenges through targeted design choices described
in Section~\ref{sec:method}, and quantify each contribution through
ablation on a held-out meta-validation set (Section~\ref{sec:exp}).

\section{RELATED WORK}
\label{sec:rw}

\textbf{SAR object detection.} Deep-learning SAR detection has been
studied extensively for ships \cite{li2017ssdd,wang2019sarship} and
aircraft \cite{chen2016sartarget}. These targets are compact (aspect
ratios $\approx$1--4:1); RFI stripes are qualitatively different. The
nearest analogous tasks are lane-line and wire/cable detection in RGB
imagery, both of which favour architectures with elongated receptive
fields.

\textbf{Small object detection.} Standard detectors struggle when
objects occupy fewer than 32$\times$32 pixels owing to limited feature
map resolution and gradient imbalance \cite{liu2021smallsurvey}.
Proposals for improvement include feature pyramid networks
\cite{lin2017fpn}, transformer-based decoders \cite{carion2020detr},
and modified regression losses. Normalised Wasserstein Distance (NWD)
\cite{wang2021nwd} replaces IoU regression with a Wasserstein distance
between Gaussian distributions centred on predicted and target boxes,
yielding non-zero gradients for non-overlapping small targets.

\textbf{Pseudo-label distillation.} Knowledge distillation via
pseudo-labels \cite{hinton2015distill,xie2020noisystudent} trains a
student on predictions from a stronger teacher, transferring soft
knowledge encoded in the teacher's confidence scores. When the teacher
is a multi-fold ensemble and the student is fine-tuned on unlabelled
in-distribution data, the approach is sometimes called noisy student
training \cite{xie2020noisystudent} or self-training
\cite{zoph2020rethinking}.

\textbf{Test-time augmentation and WBF.} Ensemble methods aggregate
predictions across multiple checkpoints or image transformations.
WBF \cite{solovyev2021wbf} clusters overlapping detections and computes
a weighted average of box coordinates, avoiding the precision-recall
trade-off of NMS-based fusion, and is widely used when combining
predictions at multiple resolutions.

\section{METHOD}
\label{sec:method}

The pipeline has four stages: base training, pseudo-label generation,
fine-tuning with distillation, and multi-resolution TTA ensemble. It is
summarised in Figure~\ref{fig:pipeline}.

\begin{figure*}[t]
\centering
\setlength{\fboxsep}{3pt}
\scriptsize
\begin{tabular}{@{}c@{$\,\to\,$}c@{$\,\to\,$}c@{$\,\to\,$}c@{$\,\to\,$}c@{$\,\to\,$}c@{$\,\to\,$}c@{}}
\fbox{\parbox[c][12mm][c]{17mm}{\centering YOLO26x base\\5 folds\\CLAHE\,+\,NWD}} &
\fbox{\parbox[c][12mm][c]{15mm}{\centering V12 teacher\\WBF\\ensemble}} &
\fbox{\parbox[c][12mm][c]{17mm}{\centering Pseudo-label\\786 test imgs\\conf\,$\ge$\,0.75}} &
\fbox{\parbox[c][12mm][c]{17mm}{\centering Student FT\\$f\!\in\!\{1,2,4\}$\\freeze 10}} &
\fbox{\parbox[c][12mm][c]{18mm}{\centering Multi-res TTA\\608/640/\\672/704\,+\,WBF}} &
\fbox{\parbox[c][12mm][c]{17mm}{\centering Cross-fold WBF\\drop $f_0$\\$\to S^{*}$}} &
\fbox{\parbox[c][12mm][c]{19mm}{\centering Blend $w{=}0.25$\\$S^{*}\!\oplus$\,V12\\$\to$ V17 (0.4776)}} \\
\end{tabular}
\caption{Pipeline overview. A YOLO26x base detector (CLAHE input,
pixel-space NWD loss) is trained per fold; the strongest
pre-distillation ensemble (V12) pseudo-labels the 786 test images;
per-fold students are fine-tuned on the fold-augmented data, run
multi-resolution TTA fused by WBF, fused across folds (dropping the
domain-overlapping fold~0) into a super-source $S^{*}$, and blended
with V12 at $w{=}0.25$ to give the final submission V17.}
\label{fig:pipeline}
\end{figure*}

\subsection{Data preparation}

\textbf{Dataset split.} The 3{,}154 labelled images are divided into
five stratified folds. We reserve a \emph{meta-validation set} of 401
images sampled from the fold-0 validation pool (random seed 42) for
leak-isolated local evaluation. All 2{,}753 remaining images form the
meta-training pool.

\textbf{SAR-specific augmentation.} SAR side-looking geometry imposes
that left--right reflections change the physical meaning of the image
(near-range versus far-range swaps), while vertical (azimuth) flips
change the look direction. We therefore set $\texttt{fliplr}=0$ and
$\texttt{flipud}=0$. Ablation confirmed that enabling horizontal flip
degrades mAP$_{50\text{-}95}$ by 6.5\% relative. Geometric
augmentations are restricted to conservative translations ($\pm5\%$ of
image size) and mild scale perturbations ($\pm25\%$), preventing thin
RFI boxes from shrinking below reliable detection thresholds. Colour
jitter (HSV) is suppressed because SAR quicklooks carry no chromatic
information. Base training additionally uses an offline SAR-aware
copy-paste augmentation (one augmented copy per image, three RFI
pastes per copy, max paste IoU 0.05); an ablation showed it gives no
gain on top of pseudo-distillation, so it is disabled during
fine-tuning.

\textbf{Contrast enhancement (CLAHE).} As an offline preprocessing
step we apply Contrast-Limited Adaptive Histogram Equalisation
(CLAHE, clip limit 3.0, $8\times8$ tile grid) to the luminance (L)
channel in LAB colour space, sharpening thin low-contrast stripes
against speckle without amplifying the global background.

\subsection{Base detector: YOLO26x with NWD loss}
\label{sec:base}

We use YOLO26x \cite{ultralytics2024yolo11} ($\approx$55.7\,M
parameters) pre-trained on COCO as the base detector, without the
optional P2 detection head: YOLO26's built-in small-target modules
(STAL and progressive loss) already address small objects, and an
ablation showed the added P2 head plateaued well below baseline on
this data. The risk of overfitting the 2{,}753-image training set is
addressed by the pseudo-label distillation stage
(Section~\ref{sec:pseudo}) rather than by shrinking the backbone. The
detector layout is summarised in Figure~\ref{fig:arch}.

\begin{figure*}[t]
\centering
\setlength{\fboxsep}{3pt}
\scriptsize
\begin{tabular}{@{}c@{$\,\to\,$}c@{$\,\to\,$}c@{$\,\to\,$}c@{$\,\to\,$}c@{}}
\fbox{\parbox[c][14mm][c]{19mm}{\centering Input\\$640{\times}640$\,RGB\\(CLAHE)}} &
\fbox{\parbox[c][14mm][c]{22mm}{\centering Backbone\\YOLO26 CSP-Darknet\\with C2PSA blocks}} &
\fbox{\parbox[c][14mm][c]{22mm}{\centering Neck\\PAN-style\\multi-scale fusion}} &
\fbox{\parbox[c][14mm][c]{23mm}{\centering Heads P3/P4/P5\\STAL\,+\,ProgLoss\\(no P2 head)}} &
\fbox{\parbox[c][14mm][c]{28mm}{\centering Loss\\$\mathcal{L}_{\mathrm{box}}{=}\tfrac{1}{2}\mathcal{L}_{\mathrm{CIoU}}$\\$+\tfrac{1}{2}\mathcal{L}_{\mathrm{NWD}}$,\,Cls: BCE}} \\
\end{tabular}
\caption{YOLO26x base-detector architecture used in this work: a
CSP-Darknet backbone with C2PSA attention blocks, a PAN-style
multi-scale neck, and three detection heads at strides 8/16/32 (the
optional P2 head is disabled, since YOLO26's built-in STAL and
progressive-loss modules already address small targets). The
box-regression loss blends CIoU with pixel-space Normalised
Wasserstein Distance (NWD); classification uses binary cross-entropy.}
\label{fig:arch}
\end{figure*}

\textbf{NWD regression loss.} Standard YOLO regression uses Complete
IoU (CIoU) \cite{zheng2020ciou} loss. For thin stripes where predicted
and target boxes frequently do not overlap vertically, CIoU is
identically zero, providing no gradient. We replace it with a blend of
CIoU and pixel-space NWD,
\begin{equation}
  \mathcal{L}_{\mathrm{box}} = (1-\alpha)\,\mathcal{L}_{\mathrm{CIoU}}
   + \alpha\,\mathcal{L}_{\mathrm{NWD}},\qquad \alpha=0.5.
\end{equation}
Following \cite{wang2021nwd}, each box is modelled as a 2-D Gaussian
$p(b)=\mathcal{N}(\mu_b,\Sigma_b)$ with mean $\mu_b=(c_x,c_y)^\top$ and
diagonal covariance $\Sigma_b=\mathrm{diag}(w^2/4,\,h^2/4)$, and NWD is
defined via the closed-form Wasserstein-2 distance,
\begin{equation}
  \mathcal{L}_{\mathrm{NWD}}(b,\hat b)
   = 1 - \exp\!\left(-\frac{\sqrt{W_2^2(p_b,p_{\hat b})}}{C}\right),
\end{equation}
where $C=12.8$\,pixels is a normalisation constant. An implementation
detail to note is that ultralytics' \texttt{BboxLoss.forward} receives
bounding boxes in \emph{stride-relative grid units}, not normalised
or pixel coordinates. We multiply by the anchor stride before
computing NWD so that $C$ is interpreted in pixel space, which keeps
gradients non-saturating in the small-height regime characteristic of
RFI stripes (median $\approx$10\,px).

\textbf{Hyperparameters.} Base training runs for up to 120 epochs
(early stopping, patience 20) with the MuSGD optimiser, initial
learning rate $\eta_0{=}10^{-3}$, cosine decay to $\eta_f{=}10^{-5}$,
3 warmup epochs, batch 16, 640\,px resolution, box-loss weight 12.0.

\subsection{Pseudo-label distillation}
\label{sec:pseudo}

\textbf{Teacher ensemble.} The teacher (denoted V12) is our strongest
pre-distillation ensemble: a single-stage WBF over
multiple YOLO26x fold and multi-resolution-TTA sources (with the
fold-2 source replaced by a CLAHE-retrained variant) together with one
RF-DETR-Large \cite{roboflow2024rfdetr} fold checkpoint weighted
$1.5\times$,
\begin{equation}
  \hat D_{\mathrm{teacher}}
  = \mathrm{WBF}\!\left(\{D_f\}_{f},\,D_{\mathrm{RF}};\;\tau{=}0.70,\,
    \texttt{conf\_type}{=}\texttt{max}\right),
\end{equation}
where $\tau$ is the IoU clustering threshold. V12 scores 0.4755 on the
held-out test set. WBF parameters ($\tau{=}0.70$,
\texttt{conf\_type=max}) were fixed after sweeping
$\tau\!\in\!\{0.65,0.70,0.75\}$ and four confidence aggregation modes;
all variants outside $(0.70,\,\texttt{max})$ degraded performance. The
marginal value of the RF-DETR component is ablated in
Section~\ref{sec:exp}.

\textbf{Pseudo-labelling the test set.} The teacher is applied to the
786 unlabelled test images with a confidence threshold of 0.75,
keeping at most 20 boxes per image, producing 600 pseudo-labelled
images after discarding low-density detections. These are merged with
each fold's training split to form \emph{fold-augmented} datasets.

\textbf{Student fine-tuning.} For each fold $f\in\{1,2,4\}$, a student
is fine-tuned from the per-fold base checkpoint of
Section~\ref{sec:base} on the fold-augmented dataset, with the
backbone frozen for the first 10 epochs (\texttt{freeze=10}), AdamW \cite{loshchilov2019adamw}
$\eta_0{=}10^{-4}$, 15 total epochs, batch 24, and the same NWD loss
blend as base training. Fold 0 is omitted because its base partition
overlaps the meta-validation pool: the fold-0 fine-tuned student
produces correlated rather than complementary predictions, and
including it empirically degrades ensemble accuracy
(Section~\ref{sec:exp}).

\subsection{Multi-resolution TTA and final fusion}

Each fine-tuned student is applied to the meta-validation and test
images at four resolutions $\mathcal{I}=\{608,640,672,704\}$ pixels.
Predictions across the four resolutions are fused per-fold via WBF
(IoU=0.70, \texttt{conf\_type=max}), producing one TTA-fused
prediction set per fold. Resolutions above 704\,px were tested and
found to degrade performance due to amplified speckle aliasing at
YOLO's feature pyramid strides. Flip-based TTA (horizontal/vertical)
is explicitly disabled.

Per-fold TTA predictions are then fused into a single super-source
$S^*$ via WBF with equal weights,
\begin{equation}
  S^* = \mathrm{WBF}\!\left(D^{\mathrm{TTA}}_1, D^{\mathrm{TTA}}_2,
   D^{\mathrm{TTA}}_4;\,w{=}1.0,\,\tau{=}0.70\right).
\end{equation}
The final submission candidate (denoted V17) is a two-source blend of
the original teacher output and the super-source,
\begin{equation}
  D_{\mathrm{final}} = \mathrm{WBF}\!\left(D_{\mathrm{V12}},\,S^*;
   \,w{=}(1.0,\,0.25),\,\tau{=}0.70\right).
\end{equation}
The blend weight $w{=}0.25$ on the super-source was selected by
sweeping $w\!\in\!\{0.10,0.15,0.25,0.35,0.50\}$ on the
meta-validation set only.

\section{EXPERIMENTS}
\label{sec:exp}

\subsection{Evaluation protocol}

\textbf{Metric.} COCO mAP$_{50\text{-}95}$ computed with pycocotools
\cite{lin2014coco} on the 401-image meta-validation set. Held-out
test-set mAP is measured by the challenge platform on the undisclosed
test partition: submissions are uploaded as a zipped COCO-format JSON
of detections, and the platform returns the official score.

\textbf{Leak isolation.} Cross-fold validation mAPs on the full
training set are inflated (``leaked'') because each fold model has seen
fold-complementary images during training. The 401-image meta-val is
honest only for models whose training partition does not overlap it.
Rows in Table~\ref{tab:ablation} for which the model has trained on
images sampled into the meta-validation pool are therefore subject to
partial leakage; this is consistent with the observed inversion at the
last row, where the leaked teacher ceiling caps reported meta-val
mAP$_{50\text{-}95}$.

\subsection{Component ablation}

\begin{table}[t]
\centering
\caption{Component ablation on 401-image meta-validation.}
\label{tab:ablation}
\small
\begin{tabular}{lcc}
\toprule
\textbf{Configuration} & \textbf{mAP$_{50\text{-}95}$} & $\Delta$ \\
\midrule
Cold fold-0 baseline                  & 0.394 & --- \\
+ Pseudo-FT, fold 4 only              & 0.435 & $+0.041$ \\
+ Multi-res TTA, fold 4 only          & 0.469 & $+0.034$ \\
+ 3-fold (drop F0), equal weights     & 0.471 & $+0.002$ \\
+ V12 blend @ $w{=}0.25$ (V17)        & 0.471$^\dagger$ & --- \\
\bottomrule
\end{tabular}\\[2pt]
{\scriptsize $^\dagger$Meta-val cannot exceed V12 by adding an honest
source (V12 is leaked).\\
Official held-out test: \textbf{0.4776}.}
\end{table}

Table~\ref{tab:ablation} reports each component addition. The chain
$0.394 \to 0.435 \to 0.469 \to 0.471$ totals $+0.077$ over cold
training; the V12 blend does not improve meta-val (the teacher
already saturates the leaked ceiling) but it transfers to the
held-out test set, where the same blend obtains 0.4776 versus 0.4755
for V12 alone.

\subsection{Fold selection ablation}

\begin{table}[t]
\centering
\caption{Cross-fold ensemble: effect of fold-0 inclusion.}
\label{tab:folds}
\small
\begin{tabular}{lc}
\toprule
\textbf{Fold set} & \textbf{mAP$_{50\text{-}95}$} \\
\midrule
Fold 4 alone (TTA)                & 0.469 \\
4-fold: F0+F1+F2+F4 equal weights & 0.458 \\
3-fold: F1+F2+F4 (drop F0)        & \textbf{0.471} \\
\bottomrule
\end{tabular}
\end{table}

Table~\ref{tab:folds} shows that the 4-fold ensemble is \emph{worse}
than fold-4 alone, and the 3-fold (drop-F0) is best. This is
explained by the domain overlap of fold-0 with the meta-validation
pool: the fold-0 fine-tuned student produces correlated, low-diversity
predictions that drag down the consensus.

\subsection{Blend weight sweep}

\begin{table}[t]
\centering
\caption{Blend weight sweep: V12 @ $1.0$ + super-source @ $w$.}
\label{tab:blend}
\small
\begin{tabular}{ccc}
\toprule
\textbf{$w$} & \textbf{Leaked fold-0 val} & \textbf{Official test} \\
\midrule
0.00 (V12 alone) & 0.5154 & 0.4755 \\
0.25             & 0.5148 & \textbf{0.4776} \\
0.375            & 0.5135 & 0.4768 \\
\bottomrule
\end{tabular}
\end{table}

Table~\ref{tab:blend} reports the leaked fold-0 validation
mAP$_{50\text{-}95}$ at three super-source blend weights together with
the corresponding official held-out test scores. The V12 source is
itself leaked on fold-0 val, so the baseline at $w{=}0$ is the leaked
ceiling; adding the honest super-source decreases leaked val
monotonically while improving held-out test up to $w{=}0.25$.

\subsection{Negative results}

We systematically evaluated several axes that did not transfer to the
held-out test set.

\textbf{RF-DETR ensemble.} The teacher's RF-DETR-Large
\cite{roboflow2024rfdetr} component (fine-tuned from COCO, weight
$1.5\times$) improves leaked fold-0 validation mAP by only $+0.001$
over the YOLO-only ensemble, and the official test set by $+0.0002$;
across six checkpoint--weight variants the RF-DETR axis is saturated,
so it is retained at minimal weight but contributes negligibly.

\textbf{D-FINE-X.} D-FINE-X \cite{peng2024dfine}
(Obj365$\to$COCO pre-training) was fine-tuned on fold 0 for ten epochs.
Final mAP$_{50\text{-}95}$ reached $3{\times}10^{-4}$, more than
$1{,}000\times$ below the fold-0 baseline, despite a clean loss
decrease from 60 to 30. DETR-family decoders appear to require many
more epochs to adapt to the elongated aspect-ratio geometry of RFI
boxes.

\textbf{Score calibration.} Temperature-scaled post-hoc calibration
applied to the final blend yielded at most $+0.0002$ mAP, consistent
with the teacher ensemble already encoding well-calibrated confidence
scores.

\textbf{Inference-only SAHI.} Sliced-inference SAHI
\cite{akyon2022sahi} at 320-pixel tiles regressed mAP by $-0.085$,
consistent with tile-boundary fragmentation of long horizontal
stripes.

\section{DISCUSSION AND CONCLUSION}

The honest per-fold baseline saturates around 0.394
mAP$_{50\text{-}95}$ (Table~\ref{tab:ablation}, top row): 2{,}753
images are too few to support reliable localisation of thin,
low-contrast RFI stripes across the full IoU=0.50--0.95 sweep. Three
findings drove the eventual lift to 0.4776. \emph{(i)}
Pseudo-distillation from a multi-source ensemble teacher contributed
the largest single gain ($+0.041$): the soft signal in the teacher's
confidence scores compensates for the small training set more
effectively than the architectural alternatives we evaluated.
\emph{(ii)}
Multi-resolution TTA between 608 and 704\,px added a further $+0.034$
because thin stripes a few pixels tall are resolved unevenly by YOLO's
strided feature pyramid at any single input size. \emph{(iii)}
Dropping the domain-overlapping fold-0 from the cross-fold ensemble
--- counter-intuitively, a \emph{smaller} ensemble --- was
net-positive, because correlated predictions drag the WBF consensus
when one source has seen the validation data during training.

Several axes that look attractive on leaked cross-fold validation did
not transfer to held-out test: a second architectural family (RF-DETR)
plateaued at $+0.0002$ across six checkpoint--weight variants;
inference-only SAHI
fragmented long horizontal stripes at tile boundaries ($-0.085$); and
temperature-scaled calibration could not improve a teacher whose
confidences were already well-conditioned by WBF fusion.
Horizontal-flip TTA cost $-6.5\%$ relative, consistent with SAR
side-looking geometry being asymmetric in azimuth--range
\cite{oliver2004sar}. The largest remaining gap is at high IoU
thresholds, where sub-pixel height accuracy dominates and a 640-px
input under-resolves the median $\approx$10\,px stripe; ranked by
expected return, the most promising next moves are (a) sliced
\emph{training} at 320-px tiles so the backbone sees RFI at doubled
relative height, (b) a DETR-family detector with a different decoder
trained for substantially more epochs than the ten that proved
insufficient for D-FINE-X \cite{huang2024deim}, and (c) a second
pseudo-label refresh with V17 itself as the new teacher.

The final pipeline --- YOLO26x with pixel-space NWD regression loss,
ensemble-teacher pseudo-distillation, multi-resolution TTA and a
selective cross-fold WBF --- reaches 0.4776 mAP$_{50\text{-}95}$ on
the held-out test set; on meta-validation the ablation chain in
Table~\ref{tab:ablation} totals $+0.077$ over the per-fold baseline.
The submission placed the \emph{UdeM AI} team 28th of 120 teams on
the public leaderboard of ESA ClearSAR Track~1. All training and fine-tuning were performed on
a single cloud NVIDIA A100 GPU.

\vfill\pagebreak

\bibliographystyle{IEEEbib}
\begin{thebibliography}{99}

\bibitem{njoku2005amsre}
E.~G. Njoku, P.~Ashcroft, T.~K. Chan, and L.~Li,
``Global survey and statistics of radio-frequency interference in
AMSR-E land observations,''
\emph{IEEE Trans. Geosci. Remote Sens.}, vol.~43, no.~5,
pp.~938--947, 2005.

\bibitem{clearsar2026}
European Space Agency,
``ClearSAR Track 1: Radio-frequency interference detection in
Sentinel-1 quicklook imagery,''
Challenge specification, 2026.

\bibitem{lin2014coco}
T.-Y. Lin, M.~Maire, S.~Belongie, L.~Bourdev, R.~Girshick, J.~Hays,
P.~Perona, D.~Ramanan, C.~L. Zitnick, and P.~Doll\'ar,
``Microsoft COCO: Common objects in context,''
in \emph{Proc. ECCV}, 2014, pp.~740--755.

\bibitem{li2017ssdd}
J.~Li, C.~Qu, and J.~Shao,
``Ship detection in SAR images based on an improved Faster R-CNN,''
in \emph{Proc. SAR Big Data Era (BIGSARDATA)}, 2017.

\bibitem{wang2019sarship}
Y.~Wang, C.~Wang, H.~Zhang, Y.~Dong, and S.~Wei,
``A SAR dataset of ship detection for deep learning under complex
backgrounds,''
\emph{Remote Sensing}, vol.~11, no.~7, p.~765, 2019.

\bibitem{chen2016sartarget}
S.~Chen, H.~Wang, F.~Xu, and Y.-Q. Jin,
``Target classification using the deep convolutional networks for SAR
images,''
\emph{IEEE Trans. Geosci. Remote Sens.}, vol.~54, no.~8,
pp.~4806--4817, 2016.

\bibitem{liu2021smallsurvey}
Y.~Liu, P.~Sun, N.~Wergeles, and Y.~Shang,
``A survey and performance evaluation of deep learning methods for
small object detection,''
\emph{Expert Systems with Applications}, vol.~172, p.~114602, 2021.

\bibitem{lin2017fpn}
T.-Y. Lin, P.~Doll\'ar, R.~Girshick, K.~He, B.~Hariharan, and
S.~Belongie,
``Feature pyramid networks for object detection,''
in \emph{Proc. CVPR}, 2017, pp.~2117--2125.

\bibitem{carion2020detr}
N.~Carion, F.~Massa, G.~Synnaeve, N.~Usunier, A.~Kirillov, and
S.~Zagoruyko,
``End-to-end object detection with transformers,''
in \emph{Proc. ECCV}, 2020, pp.~213--229.

\bibitem{wang2021nwd}
J.~Wang, C.~Xu, W.~Yang, and L.~Yu,
``A normalized Gaussian Wasserstein distance for tiny object
detection,''
\emph{arXiv:2110.13389}, 2021.

\bibitem{hinton2015distill}
G.~Hinton, O.~Vinyals, and J.~Dean,
``Distilling the knowledge in a neural network,''
\emph{arXiv:1503.02531}, 2015.

\bibitem{xie2020noisystudent}
Q.~Xie, M.-T. Luong, E.~Hovy, and Q.~V. Le,
``Self-training with noisy student improves ImageNet classification,''
in \emph{Proc. CVPR}, 2020, pp.~10687--10698.

\bibitem{zoph2020rethinking}
B.~Zoph, G.~Ghiasi, T.-Y. Lin, Y.~Cui, H.~Liu, E.~D. Cubuk, and
Q.~V. Le,
``Rethinking pre-training and self-training,''
in \emph{Proc. NeurIPS}, 2020.

\bibitem{solovyev2021wbf}
R.~Solovyev, W.~Wang, and T.~Gabruseva,
``Weighted boxes fusion: Ensembling boxes from different object
detection models,''
\emph{Image and Vision Computing}, vol.~107, p.~104117, 2021.

\bibitem{ultralytics2024yolo11}
G.~Jocher and J.~Qiu,
``Ultralytics YOLO26,''
\url{https://github.com/ultralytics/ultralytics}, 2026.

\bibitem{zheng2020ciou}
Z.~Zheng, P.~Wang, W.~Liu, J.~Li, R.~Ye, and D.~Ren,
``Distance-IoU loss: Faster and better learning for bounding box
regression,''
in \emph{Proc. AAAI}, 2020, pp.~12993--13000.

\bibitem{roboflow2024rfdetr}
Roboflow,
``RF-DETR: A real-time transformer-based object detector,''
\url{https://github.com/roboflow/rf-detr}, 2024.

\bibitem{loshchilov2019adamw}
I.~Loshchilov and F.~Hutter,
``Decoupled weight decay regularization,''
in \emph{Proc. ICLR}, 2019.

\bibitem{peng2024dfine}
Y.~Peng, H.~Li, P.~Wu, Y.~Zhang, X.~Sun, and F.~Wu,
``D-FINE: Redefine regression task of DETR as fine-grained
distribution refinement,''
\emph{arXiv:2410.13842}, 2024.

\bibitem{akyon2022sahi}
F.~C. Akyon, S.~O. Altinuc, and A.~Temizel,
``Slicing aided hyper inference and fine-tuning for small object
detection,''
in \emph{Proc. IEEE ICIP}, 2022, pp.~966--970.

\bibitem{oliver2004sar}
C.~Oliver and S.~Quegan,
\emph{Understanding Synthetic Aperture Radar Images}.
SciTech Publishing, 2004.

\bibitem{huang2024deim}
S.~Huang, Z.~Lu, X.~Cun, Y.~Yu, X.~Zhou, and X.~Shen,
``DEIM: DETR with improved matching for fast convergence,''
\emph{arXiv:2412.04234}, 2024.

\end{thebibliography}

\end{document}
